% ============================================================
% BÖLÜM 5 - SONUÇLAR VE TARTIŞMA
% ============================================================
\chapter{SONUÇLAR VE TARTIŞMA}
\label{ch:sonuc}

\section{Genel Değerlendirme}

Bu tez kapsamında, fiziksel operatörlerin sinir ağı yapısına doğrudan gömülmesi prensibi üzerine inşa edilen INNATE fiziksel-operatör tabanlı sinir operatörü, ısı transferi içeren çok ölçekli bir karışık konveksiyon probleminde geliştirilmiş, eğitilmiş ve bağımsız bir LES referans veri seti ile doğrulanmıştır. Üretilen model 9905 öğrenilebilir parametre içermekte olup 20 katmanlı fraksiyonel adım IMEX zaman integratörü üzerinde işlemektedir. Kademe~1+2+3 disiplinli eğitim stratejisi sayesinde, ön deneylerde gözlemlenen anti-fiziksel davranışlar (parametre kaçışı, Nusselt sayısı abartımı, genlik pompalama) sistematik biçimde giderilmiş ve model, LES referansı ile \emph{nitel olarak tutarlı} (enerji spektrumu monoton-azalan; ısı transferi yönü doğru; uzun zaman ufkunda kararlı diverjans-serbest çözüm üreten) ancak \emph{niceliksel olarak yaklaşık $2\times$ (TKE, Nu) ve $4$-$5\times$ ($\theta_{\mathrm{rms}}$) sapma gösteren} bir noktaya getirilmiştir.

Bu sonuçlar, üç temel iddiayı destekler nitelikte değerlendirilebilir.

\paragraph{İddia 1 --- Fiziksel-operatör tabanlı mimari karışık konveksiyon problemlerinde uygulanabilirdir.} Fiziksel operatörlerin yapısal olarak temsil edilmesi (mimari açıdan saf-fizik; eğitim açısından bağımsız LES referansı ile destekli), modelin diverjanssızlık ve sayısal kararlılık gibi temel akışkanlar mekaniği garantilerini otomatik olarak korumasını sağlamıştır. Eğitim sürecinde hiçbir NaN/Inf koruma tetiklenmemiş, diverjans hatası LES referansının sayısal toleransı içinde kalmıştır. Bu, klasik PINN paradigmalarında uzun zaman ufuklarında sıkça gözlenen ``yumuşak kısıt biriken hatası'' probleminin yapısal düzeyde çözülmüş olduğunu gösterir.

\paragraph{İddia 2 --- Disiplinli eğitim, mimari kadar belirleyicidir.} Kademe~1 parametre hijyeni olmaksızın, model 120 öğrenilebilir parametreyi kullanarak kanonik Boussinesq--Newton denklem yapısından sapacak şekilde kayıp fonksiyonunu \emph{aldatmıştır}. Kademe~2 PDE artık kaybı olmaksızın, model alt-ızgara modelleme katsayılarını fiziksel olarak makul aralıkta tutmakta zorlanmıştır. Kademe~3 minimal 5-terimli kayıp seti olmaksızın, optimizer en büyük gradyanı üreten kayıp tarafından sürüklenmiş ve ikincil fizik kısıtları görmezden gelinmiştir. Üç kademenin birlikte uygulanması, bu üç başarısızlık modunun da ortadan kaldırılmasını sağlamıştır.

\paragraph{İddia 3 --- Düşük parametre sayılı bir model ile kararlı çözüm üretilebilir.} 9905 öğrenilebilir parametre, klasik saf veri-tabanlı operatör öğrenme modellerinin ($10^{6}$--$10^{7}$ parametre mertebesinde) parametre sayılarına kıyasla yaklaşık iki büyüklük mertebesi daha küçük bir model demektir. Bu model ile diverjans-serbest, sayısal olarak kararlı, uzun-zaman rollout sırasında patlamayan bir çözüm üretilmiştir. Niceliksel doğruluk (TKE, Nu, $\theta_{\mathrm{rms}}$ açısından) henüz mühendislik tahmini düzeyinde değildir; ancak bu sapmaların kök sebepleri (Bölüm~\ref{ch:bulgular}, §\ref{sec:termal_sinir}) somut olarak tespit edilmiş olup parametre sayısı artırılmadan da iyileştirilebilirliği gelecek çalışmaların kapsamına alınmıştır.

\section{Başarılı Yönler}

Sayısal bulgular ışığında, INNATE saf-fizik sinir operatörünün başarılı olduğu yönler şunlardır:

\begin{itemize}
  \item \textbf{Yapısal fizik garantileri.} Diverjanssızlık her zaman adımı sonunda yapısal olarak sağlanmıştır; bu, klasik PINN'lerde elde edilemeyen bir özelliktir. Diverjans hatası LES referans seviyesinde kalmış ve uzun zaman rollout sırasında birikim göstermemiştir.

  \item \textbf{Sayısal kararlılık.} Eğitim süresince 100 epoch boyunca tek bir NaN/Inf olayı yaşanmamış; rollout sırasında modelin enerji içeriği sınırlı kalmıştır. Bu, klasik MLP-tabanlı operatör öğrenme modellerinin uzun zaman ufuklarındaki kararsızlık eğilimleriyle keskin bir tezat oluşturur.

  \item \textbf{Yorumlanabilirlik kapasitesi.} Spectral-Cs v3 parametrizasyonu, modelin öğrendiği Smagorinsky sabit alanını Fourier modları üzerinde \emph{prensip olarak} doğrudan görselleştirilebilir kılar. Bu çalışmanın eğitim bütçesi altında uzaysal değişim sınırlı kalmış olsa da (Bölüm~\ref{sec:spectral_cs_aktivasyon}), parametrizasyon klasik MLP tabanlı SGS modellemesinden yapısal bir yorumlanabilirlik avantajı sunar; daha uzun eğitim ufuklarında bu avantaj nicel olarak da gözlemlenebilir hale gelecektir.

  \item \textbf{Düşük parametre sayısı.} 9905 öğrenilebilir parametre, modern operatör öğrenme modellerinden iki ile dört büyüklük mertebesi daha küçüktür. Bu, eğitim süresinin (RTX~4090 üzerinde $\sim$6 saat) ve eğitim sırasında oluşan bellek izinin sınırlı kalmasını sağlar.

  \item \textbf{Enerji-kaskat yönünün doğru öğrenilmesi.} Spektrum eğimi Kolmogorov $-5/3$ referansından anlamlı ölçüde sapma göstermekle birlikte (model: $-2.5$ ila $-3.0$; LES: $\approx -1.75$), monotonca negatif ve düzgün biçimde azalan bir enerji spektrumu üretilmiştir. Bu, modelin enerji-kaskatın yönünü doğru öğrendiğini fakat küçük ölçeklerde fazla yumuşatma uyguladığını gösterir.
\end{itemize}

\section{Sınırlamalar ve Yapısal Kısıtlar}

Akademik dürüstlük çerçevesinde, çalışmanın sınırlarını da açık biçimde belirtmek gerekmektedir.

\subsection{Termal Alanın Tam-Eşlenik Olmaması}

Bölüm~\ref{sec:termal_sinir}'te ayrıntılı olarak tartışıldığı üzere, INNATE mimarisinde sıcaklık alanı momentum alanı ile yalnızca kaldırma kuvveti terimi üzerinden ve $\mathrm{Ri}\approx 10^{-3}$ gibi zayıf bir bağlantı katsayısıyla eşleniktir. Bu yapısal sınır, $\theta_{\mathrm{rms}}$ ve $v\theta$ akısı gibi termal niceliklerin LES referansından $4$--$14$ kat sapma göstermesinin temel kaynağıdır.

Bu sınır, çalışmanın eğitim aşaması tamamlandıktan sonra, geriye dönük analizler sırasında keşfedilmiş olduğu için bu tezin kapsamı içinde mimari düzeyde çözülememiştir. Bu, çalışmanın bir başarısızlığı değil, bir \emph{öğrenmesidir}: lisans bitirme projesinde keşfedilen yapısal bir kısıtın sonraki çalışmalara aktarılması olarak değerlendirilmelidir.

\subsection{Çalışma Noktasının Tekilliği}

Çalışmanın eğitim ve doğrulama aşamaları, yalnızca iki Reynolds sayısı ($\mathrm{Re}\in\{7000,10\,000\}$) ile sınırlandırılmıştır. Bu, hesaplama kaynaklarının verimli kullanımı amacıyla bilinçli bir karardır; ancak modelin daha düşük veya daha yüksek Reynolds sayılarındaki davranışı bu çalışma kapsamında değerlendirilmemiştir. Spectral-Cs parametrizasyonunun farklı Reynolds sayılarına \emph{ekstrapolasyon} kapasitesi, gelecek çalışmaların ele alması gereken bir konudur.

\subsection{Geometri Sınırı}

Çalışma, üç boyutlu periyodik küp benzeri bir geometride yapılmıştır. Pek çok mühendislik uygulaması (örneğin kapalı kutu Rayleigh--B\'enard konveksiyonu, kanat profili çevresi akış, kanal akışı), duvar sınırlı ve periyodik olmayan geometrileri içerir. Spektral projeksiyon nöronu ve Spectral-Cs parametrizasyonu, bu tür geometrilere doğrudan uygulanabilir değildir ve önemli mimari değişiklikler gerektirir.

\subsection{Zorlama Amplitüdünün Serbest Bırakılması}

Final eğitim konfigürasyonunda, zorlama amplitüdü $A_{F}$ Kademe~1 freeze politikasının dışında bırakılmıştır. Bu kararın gerekçesi, $A_{F}$'nin denklem-yapısı parametresi değil, kaynak-büyüklüğü parametresi olmasıdır; ancak final modelin TKE değerlerinin LES'in $\sim 2$ katı olması, optimizer'in bu parametreyi yüksek değerlere ittiği şüphesini doğurmaktadır. Gelecek çalışmalarda $A_{F}$'nin de Kademe~1 disiplinine alınması, ya da alternatif olarak doğrudan TKE eşleştirme kaybının kayıp fonksiyonuna eklenmesi makul bir adım olacaktır.

\subsection{Kıyaslama Eksikliği}

Bu çalışmada INNATE modeli, yalnızca LES referansı ile karşılaştırılmıştır. Diğer makine öğrenmesi paradigmaları (Fourier neural operator, DeepONet, PINN, hibrit yöntemler) ile doğrudan sayısal karşılaştırma yapılmamıştır. Bu, çalışmanın derinlik-genişlik takasının bilinçli bir tercihidir; ancak gelecek çalışmaların, INNATE'in modern operatör öğrenme modelleriyle yan yana karşılaştırılmasını sağlaması arzu edilir.

\section{Gelecek Çalışma Önerileri}

Yukarıdaki sınırlamalar, doğrudan gelecek çalışmalara yönelik somut araştırma soruları üretir:

\paragraph{Tam-eşlenik termal mimari.} Termal alanın momentum alanı ile yapısal düzeyde \emph{simetrik} biçimde temsil edildiği bir INNATE varyantı geliştirilebilir. Somut olarak, hız alanına uygulanan adveksiyon ve projeksiyon nöronlarının paralellerinin sıcaklık alanına da uygulanması ve termal-momentum eşleniğinin yalnızca kaldırma terimi üzerinden değil, fakat tam denklem yapısı boyunca korunması incelenebilir.

\paragraph{Reynolds genelleme.} Spectral-Cs parametrizasyonunun farklı Reynolds sayılarına genelleme kapasitesi, $\mathrm{Re}\in\{3000, 5000, 15\,000, 20\,000\}$ gibi ek çalışma noktaları üzerinden sistematik biçimde sınanabilir. Eğer Spectral-Cs katsayıları $\mathrm{Re}$'ye duyarlı biçimde otomatik adaptasyon gösteriyorsa, modelin tek bir eğitim koşusu ile birden çok Reynolds sayısı bandına uygulanabilirliği değerlendirilebilir.

\paragraph{Periyodik olmayan geometri.} Spektral projeksiyon nöronunun yerini, sonlu fark veya sonlu eleman tabanlı bir alternatifin alabileceği bir INNATE varyantı geliştirilebilir; bu, modelin duvar-sınırlı geometrilere uygulanabilirliğini açar.

\paragraph{Karşılaştırmalı değerlendirme çerçevesi.} INNATE'in modern operatör öğrenme modelleri (FNO, DeepONet, hibrit yöntemler) ile aynı problem üzerinde yan yana karşılaştırılması, saf-fizik mimari prensibinin görece avantajlarının ve dezavantajlarının net biçimde ortaya konmasını sağlar.

\paragraph{Zorlama amplitüdünün disipline alınması.} $A_{F}$ parametresinin Kademe~1 freeze politikasına dahil edilmesi veya kayıp fonksiyonuna doğrudan TKE eşleştirme teriminin eklenmesi, TKE abartımının giderilmesine yönelik tek başına basit ve etkili bir adım olabilir.

\paragraph{Daha uzun zaman ufukları.} Bu çalışmadaki rollout süresi $\sim 200$ birim zaman ile sınırlıdır. Daha uzun rollout sürelerinde modelin enerji bütçesinin uzun-vadeli istikrarı incelenebilir.

\paragraph{Spectral-Cs uzaysal kapasitenin tam aktive edilmesi.} Bu çalışmanın eğitim bütçesi (RTX~4090 dizüstü GPU üzerinde 100 epoch) altında, Spectral-Cs Fourier katsayılarının büyük çoğunluğu başlangıç değerlerinden anlamlı biçimde uzaklaşmamış; modelin uzaysal değişen $C_{s}(\mathbf{x})$ kapasitesi tam olarak aktive olmamıştır (Bölüm~\ref{sec:spectral_cs_aktivasyon}). Daha güçlü donanım (örneğin çoklu-GPU veya HPC kümesi) ve daha uzun eğitim süreleri (örneğin 500--1000 epoch) altında, Fourier modlarının fiziksel olarak anlamlı uzaysal yapılar üretmesi beklenmektedir. Aynı zamanda, Fourier katsayıları için ayrı bir öğrenme hızı zamanlayıcısı veya warm-up sonrası amplifikasyon disiplini de bu kapasitenin daha hızlı aktive olmasını sağlayabilir.

\section{Kapanış}

Lisans bitirme projesinin bu ikinci dönemi, ilk dönemde TGV3D problemi üzerinden kavramsal düzeyde önerilen INNATE saf-fizik sinir operatörü yaklaşımının, ısı transferi ve çok ölçekli türbülans içeren bir karışık konveksiyon probleminde fiilen uygulanabileceğini göstermiştir. Bu süreçte saf-fizik mimari prensibinin tek başına yeterli olmadığı, eğitim disiplininin en az mimari kadar belirleyici olduğu somut olarak gözlemlenmiştir. Kademe~1+2+3 yaklaşımı, bu öğrenme sürecinin metodolojik bir kazanımıdır.

Çalışmanın akademik açıdan en önemli katkısı, fiziksel olarak yorumlanabilir, parametre-verimli ve sayısal olarak kararlı bir sinir operatörünün ısı transferi içeren karışık konveksiyon problemlerinde uygulanabilir olduğunu somut metriklerle göstermesidir. Çalışmanın açıkça belgelediği sınırlamalar (özellikle termal alanın tam-eşlenik olmaması) ise, bu yaklaşımın geliştirilmesinde bundan sonra atılması gereken adımların somut bir yol haritasını sunar.

Sonuç olarak, bu tez, fiziksel operatörlerin sinir ağı yapısına doğrudan gömülmesi yaklaşımının türbülanslı akış modellemesinde nasıl somut bir araştırma çizgisi oluşturabileceğine dair bir örnek sunmaktadır. Yaklaşımın hem başarılı olduğu yönler hem de keşfedilen yapısal kısıtlar, hesaplamalı akışkanlar mekaniği ile yapay zekâ arasındaki köprünün daha sağlam temellere oturtulması açısından değerlendirilmelidir.
